\chapter{Los bloques de construcción}
\label{cap:bloques}

\begin{enpocas}
Las seis piezas con las que se arma YOLO26, cada una como una clase de PyTorch:
\cod{Conv}, \cod{Cuello}, \cod{C3k2}, \cod{SPPF}, \cod{Atencion} y \cod{C2PSA}.
Todo está en \cod{red_fuego/bloques.py}.
\end{enpocas}

\section{Cómo se escribe una capa en PyTorch}

En PyTorch toda capa es una clase que hereda de \cod{nn.Module}. En
\cod{__init__} se crean las sub-capas (las que tienen pesos) y en
\cod{forward} se dice qué se hace con la entrada. Los datos viajan como
\textbf{tensores} de cuatro dimensiones $(B, C, H, W)$: $B$ imágenes del lote,
$C$ canales y $H\times W$ píxeles. Una imagen RGB de $640\times640$ entra como
$(1, 3, 640, 640)$.

\section{Conv: la unidad mínima}

\paso{4} \textbf{Escribir \cod{Conv}.}
Toda convolución de YOLO va seguida de una normalización por lotes
(BatchNorm) y de la activación SiLU, $\mathrm{SiLU}(x) = x\,\sigma(x)$:

\begin{lstlisting}[style=py]
class Conv(nn.Module):
    def __init__(self, c1, c2, k=1, s=1, g=1, d=1, act=True):
        super().__init__()
        self.conv = nn.Conv2d(c1, c2, k, s, relleno(k, d), groups=g, dilation=d, bias=False)
        self.bn = nn.BatchNorm2d(c2, eps=1e-3, momentum=0.03)
        self.act = nn.SiLU(inplace=True) if act else nn.Identity()

    def forward(self, x):
        return self.act(self.bn(self.conv(x)))
\end{lstlisting}

\begin{itemize}
  \item \cod{c1}, \cod{c2}: canales de entrada y de salida. \cod{k}: tamaño del filtro.
  \item \cod{s}: paso (\emph{stride}). Con \cod{s=2} el mapa sale a la mitad de alto
        y de ancho: así es como la red «se aleja» de la imagen.
  \item \cod{relleno(k)} $=\lfloor k/2\rfloor$ conserva el tamaño cuando \cod{s=1}.
  \item No lleva sesgo (\cod{bias=False}) porque BatchNorm ya suma uno.
  \item \cod{g} (grupos) permite la convolución \emph{en profundidad}
        (\cod{DWConv}): cada canal tiene su propio filtro y no se mezcla con los demás.
        Cuesta $C$ veces menos y la usa la cabeza de clases.
\end{itemize}

Una \cod{Conv} $3\times3$ de 64 a 128 canales tiene $64\cdot128\cdot9 = 73\,728$
pesos más $2\cdot128$ de BatchNorm: $73\,984$ parámetros.

\section{Cuello (bottleneck) y C3k}

\paso{5} \textbf{Escribir el cuello de botella residual.}
Dos \cod{Conv} $3\times3$ y un atajo que suma la entrada:
$\mathbf{y} = \mathbf{x} + f(\mathbf{x})$. El atajo deja pasar el gradiente sin
obstáculos y permite apilar muchas capas (la idea de ResNet).

\begin{lstlisting}[style=py]
class Cuello(nn.Module):
    def __init__(self, c1, c2, atajo=True, k=(3, 3), e=0.5):
        super().__init__()
        c_ = int(c2 * e)
        self.cv1 = Conv(c1, c_, k[0])
        self.cv2 = Conv(c_, c2, k[1])
        self.suma = atajo and c1 == c2

    def forward(self, x):
        y = self.cv2(self.cv1(x))
        return x + y if self.suma else y
\end{lstlisting}

\cod{C3k} es un CSP (\emph{Cross Stage Partial}) de tres convoluciones: la
entrada se proyecta por dos caminos; uno pasa por dos cuellos y el otro no; al
final se concatenan. Es más rico que un cuello solo y lo usan las etapas
profundas.

\section{C3k2: el bloque principal}

\paso{6} \textbf{Escribir \cod{C3k2}.}
Es el bloque que más se repite en YOLO26 (8 de las 23 capas). La
figura~\ref{fig:c3k2} muestra su flujo:

\begin{figure}[H]
\centering
\begin{tikzpicture}[font=\small,>=Stealth,
  caja/.style={draw,rounded corners=2pt,minimum height=7mm,minimum width=14mm,fill=NavyBlue!8},
  op/.style={draw,circle,inner sep=1pt,fill=Goldenrod!20}]
\node[caja] (in) {entrada $c_1$};
\node[caja,right=6mm of in] (cv1) {Conv $1{\times}1$};
\node[op,right=6mm of cv1] (sp) {$\div 2$};
\node[caja,right=8mm of sp,yshift=-9mm] (m1) {sub 1};
\node[caja,right=6mm of m1] (m2) {sub 2};
\node[right=3mm of m2] (pts) {$\cdots$};
\node[op,right=30mm of sp,yshift=9mm] (cat) {C};
\node[caja,right=6mm of cat] (cv2) {Conv $1{\times}1$};
\node[caja,right=6mm of cv2] (out) {salida $c_2$};
\draw[->] (in)--(cv1); \draw[->] (cv1)--(sp);
\draw[->] (sp) |- node[pos=0.75,above]{$c$} (cat);
\draw[->] (sp) |- node[pos=0.75,below]{$c$} (m1);
\draw[->] (m1)--(m2); \draw[->] (m2)--(pts);
\draw[->] (m1.north) -- ++(0,4mm) -| ([xshift=-1mm]cat.south);
\draw[->] (m2.north) -- ++(0,2mm) -| ([xshift=1mm]cat.south);
\draw[->] (cat)--node[above]{$(2{+}n)c$}(cv2); \draw[->] (cv2)--(out);
\end{tikzpicture}
\caption{C3k2. La entrada se proyecta a $2c$ canales y se parte en dos. La
segunda mitad pasa por $n$ subbloques en cadena; \emph{todas} las salidas
intermedias se concatenan (C) y una Conv $1{\times}1$ las mezcla. El subbloque
es un \cod{Cuello}, un \cod{C3k} o un \cod{Cuello} + atención.}
\label{fig:c3k2}
\end{figure}

\begin{lstlisting}[style=py]
class C3k2(nn.Module):
    def __init__(self, c1, c2, n=1, c3k=False, e=0.5, atencion=False, atajo=True):
        super().__init__()
        self.c = int(c2 * e)
        self.cv1 = Conv(c1, 2 * self.c, 1)
        self.cv2 = Conv((2 + n) * self.c, c2, 1)
        if atencion:
            sub = lambda: nn.Sequential(Cuello(self.c, self.c, atajo),
                                        BloquePSA(self.c, cabezas=max(self.c // 64, 1)))
        elif c3k:
            sub = lambda: C3k(self.c, self.c, 2, atajo)
        else:
            sub = lambda: Cuello(self.c, self.c, atajo)
        self.m = nn.ModuleList(sub() for _ in range(n))

    def forward(self, x):
        y = list(self.cv1(x).chunk(2, 1))      # dos mitades de c canales
        for m in self.m:
            y.append(m(y[-1]))                  # cada subbloque toma la última salida
        return self.cv2(torch.cat(y, 1))
\end{lstlisting}

\begin{consejo}
El parámetro \cod{e} (expansión) controla el ancho interno. YOLO26 usa
\cod{e=0.25} en las dos primeras etapas (mapas grandes, conviene ahorrar
cómputo) y \cod{e=0.5} en las demás.
\end{consejo}

\section{SPPF: ver a varias escalas}

\paso{7} \textbf{Escribir \cod{SPPF}.}
Al final de la columna vertebral el mapa mide $20\times20$. Para que cada celda
vea su contexto (¿esa mancha naranja está junto a una columna de humo?),
\cod{SPPF} aplica tres \emph{max-pool} de $5\times5$ en cadena. Tres
ventanas de 5 equivalen a ventanas de 5, 9 y 13 celdas, sin agregar ningún
parámetro. Las cuatro versiones se concatenan y se mezclan; YOLO26 agrega
además un atajo $\mathbf{y} \leftarrow \mathbf{y} + \mathbf{x}$.

\begin{lstlisting}[style=py]
    def forward(self, x):
        y = [self.cv1(x)]
        for _ in range(self.n):              # n = 3
            y.append(self.m(y[-1]))          # MaxPool2d(5, 1, 2)
        y = self.cv2(torch.cat(y, 1))
        return x + y if self.suma else y
\end{lstlisting}

\section{Atención: que cada celda mire a todas las demás}

\paso{8} \textbf{Escribir \cod{Atencion}, \cod{BloquePSA} y \cod{C2PSA}.}
Una convolución solo mira a sus vecinos. La \textbf{autoatención} deja que cada
celda consulte a \emph{todas} las demás. Si el mapa tiene $N = H\,W$ celdas, de
cada una se calculan tres vectores con convoluciones $1\times1$: una
\emph{consulta} $\mathbf{q}_i$, una \emph{llave} $\mathbf{k}_j$ y un
\emph{valor} $\mathbf{v}_j$. La salida de la celda $i$ es un promedio pesado de
los valores:
\begin{equation}
\mathbf{y}_i = \sum_{j=1}^{N} a_{ij}\,\mathbf{v}_j,\qquad
a_{ij} = \frac{\exp(\mathbf{q}_i\cdot\mathbf{k}_j/\sqrt{d_k})}
              {\sum_{l}\exp(\mathbf{q}_i\cdot\mathbf{k}_l/\sqrt{d_k})}.
\end{equation}
El peso $a_{ij}$ es grande cuando la consulta de $i$ «se parece» a la llave de
$j$. Con varias \emph{cabezas} se hacen varias atenciones en paralelo, cada una
con su parte de los canales. Se suma además una convolución $3\times3$ en
profundidad sobre $\mathbf{v}$, que le dice a la red dónde está cada celda
(codificación posicional).

\begin{lstlisting}[style=py]
    def forward(self, x):
        B, C, H, W = x.shape
        q, k, v = self.qkv(x).view(B, self.cabezas, 2 * self.dk + self.dc, H * W).split(
            [self.dk, self.dk, self.dc], dim=2)
        a = ((q * self.escala).transpose(-2, -1) @ k).softmax(dim=-1)   # (B, h, N, N)
        y = (v @ a.transpose(-2, -1)).view(B, C, H, W) + self.pe(v.reshape(B, C, H, W))
        return self.proy(y)
\end{lstlisting}

\begin{aviso}
La matriz $a$ tiene $N\times N$ elementos. Con el mapa de $20\times20$ son
$400^2 = 160\,000$ por cabeza: barato. Con el de $80\times80$ serían
$41$ millones. Por eso YOLO26 solo pone atención en la etapa más profunda
($P_5$).
\end{aviso}

\cod{BloquePSA} es como un bloque de transformador: atención con atajo y luego
una pequeña red de avance (dos Conv $1\times1$) también con atajo.
\cod{C2PSA} aplica la atención solo a la mitad de los canales (estilo CSP) y la
otra mitad pasa directo.

\paso{9} \textbf{Probar cada bloque.}
Antes de ensamblar nada hay que comprobar que cada bloque recibe y entrega las
formas correctas:
\begin{lstlisting}[style=py]
import torch
from cuadrirotor_sim.red_fuego.bloques import Conv, Cuello, C3k2, SPPF, C2PSA
x = torch.randn(1, 64, 80, 80)
for nombre, b in [('Conv 3x3 s=2', Conv(64, 128, 3, 2)), ('Cuello', Cuello(64, 64)),
                  ('C3k2 n=1', C3k2(64, 128, 1, False, 0.25)), ('C3k2 c3k', C3k2(64, 128, 1, True)),
                  ('SPPF', SPPF(64, 64, atajo=True)), ('C2PSA', C2PSA(64, 64))]:
    y = b(x)
    print(nombre, tuple(x.shape), '->', tuple(y.shape), sum(p.numel() for p in b.parameters()))
\end{lstlisting}

\begin{resultado}
\begin{lstlisting}
Conv 3x3 s=2   (1, 64, 80, 80) -> (1, 128, 40, 40)   73984 parámetros
Cuello         (1, 64, 80, 80) -> (1, 64, 80, 80)   37056 parámetros
C3k2 n=1       (1, 64, 80, 80) -> (1, 128, 80, 80)   26080 parámetros
C3k2 c3k       (1, 64, 80, 80) -> (1, 128, 80, 80)   78848 parámetros
SPPF           (1, 64, 80, 80) -> (1, 64, 80, 80)   10432 parámetros
C2PSA          (1, 64, 80, 80) -> (1, 64, 80, 80)   16352 parámetros
\end{lstlisting}
\end{resultado}

\begin{ejercicio}
Cambia \cod{e=0.25} por \cod{e=0.5} en el \cod{C3k2 n=1} y vuelve a contar los
parámetros. ¿Por qué crecen casi cuatro veces si solo duplicaste el ancho?
(Pista: en una convolución los pesos son $c_{\text{in}}\cdot c_{\text{out}}\cdot k^2$.)
\end{ejercicio}
